\lysh{15081_SOLUTION}{1}
\(\alpha\) Ο κύκλος \(C_1\) έχει κέντρο \(K\left(-\frac{A}{2}, \frac{B}{2}\right)\), δηλαδή \(K\left(-\sqrt{2}, 0\right)\) και ακτίνα :
\[ \rho_1 = \sqrt{\frac{A^2+B^2-4\Gamma}{2}} = 1 \]
και ο κύκλος \(C_2\) έχει κέντρο \(\Lambda\left(\frac{A}{2}, \frac{B}{2}\right)\), δηλαδή \(\Lambda\left(3\sqrt{2}, 0\right)\) και ακτίνα :
\[ \rho_2 = \sqrt{\frac{A^2+B^2-4\Gamma}{2}} = 3. \]

\(\beta\) i. Μια ευθεία που διέρχεται από την αρχή των αξόνων \(O(0,0)\) και δεν είναι κάθετη στον \(x'x\) άξονα έχει εξίσωση:
\[ (\eta): y = \lambda x \iff y - \lambda x = 0. \]
Η ευθεία \((\eta)\) εφάπτεται στους δύο κύκλους αν και μόνο αν οι αποστάσεις των κέντρων \(K\) και \(\Lambda\) από την ευθεία αυτή είναι ίσες με τις αντίστοιχες ακτίνες των κύκλων. Δηλαδή έχουμε:
\[ d(K, \eta) = 1 \quad (1) \quad \text{και} \]
\[ d(\Lambda, \eta) = 3 \quad (2). \]
Λύνουμε το σύστημα των εξισώσεων (1) και (2):
\[
\left\{
\begin{array}{l}
\left|\frac{0+\lambda\sqrt{2}}{\sqrt{1+\lambda^2}}\right| = 1 \\
\left|\frac{0-3\sqrt{2}\lambda}{\sqrt{1+\lambda^2}}\right| = 3
\end{array}
\right\}
\implies
\left\{
\begin{array}{l}
\frac{\sqrt{2}|\lambda|}{\sqrt{1+\lambda^2}} = 1 \\
\frac{3\sqrt{2}|\lambda|}{\sqrt{1+\lambda^2}} = 3
\end{array}
\right\}
\implies
\left\{
\begin{array}{l}
2\lambda^2 = 1+\lambda^2 \\
18\lambda^2 = 9+9\lambda^2
\end{array}
\right\}
\implies
\left\{
\begin{array}{l}
\lambda^2 = 1 \\
9\lambda^2 = 9
\end{array}
\right\}
\implies \lambda = \pm 1.
\]
\(`Άρα, από την αρχή των αξόνων διέρχονται δυο κοινές εφαπτόμενες των κύκλων, με εξισώσεις:`\)
\[ (\eta_1): y = -x \quad \text{και} \quad (\eta_2): y = x. \]

ii. Η αρχή των αξόνων \((0,0)\) είναι εσωτερικό σημείο της διακέντρου \(K\Lambda\), διότι η \(K\Lambda\) είναι πάνω στον άξονα \(x'x\) και έχει άκρα τα σημεία \(K(-\sqrt{2}, 0)\) και \(\Lambda(3\sqrt{2}, 0)\). Επομένως οι εφαπτόμενες που βρήκαμε στο \(\beta)\) ερώτημα είναι εσωτερικές, όπως φαίνεται στο παρακάτω σχήμα.

% TODO: TikZ σχήμα
\end{lysh}